% intro.tex -- funnel -> bottleneck -> open question -> contributions (TEACHING EXAMPLE).
\section{Introduction}\label{sec:intro}
TFHE~\cite{CGGI16,CGGI20} refreshes a ciphertext with a \emph{blind rotation} of a test
polynomial, a technique going back to FHEW~\cite{DM15}. Choosing the test polynomial evaluates
an arbitrary function during the refresh (programmable bootstrapping, \PBSop)~\cite{CJP21}.
The blind rotation dominates the cost, and it does not depend on the function: evaluating
$k$ functions of the \emph{same} input with $k$ independent \PBSop{} repeats it $k$ times.

A 4-bit S-box is the textbook instance. If the next layer is bit-sliced, the four output bits
are needed as four ciphertexts, i.e.\ four \PBSop{}. (If a single ciphertext of the nibble
suffices, one \PBSop{} with $f=S$ does the job and there is nothing to optimise.)
Carpov, Izabach\`ene and Mollimard~\cite{CIM19} observed that the test polynomials of many
functions share a common factor $v_0$, so one blind rotation of $v_0$ can serve all of them
(multi-value PBS, \MVPBS). Tree-based and multi-value variants for larger tables were studied
in~\cite{GBA21}; CMux-tree evaluation from GGSW inputs goes back to~\cite{CGGI17}.

\subsection{Contributions (of this toy study)}
\begin{itemize}
\item A self-contained derivation of the factorisation for the one-padding-bit encoding
  (\Cref{lem:factor}), the exact bound $\|d_f\|^2\le p+2$ for Boolean tables
  (\Cref{lem:norm}), and the output-variance formula (\Cref{thm:noise}).
\item A falsification-first evaluation: 3.86×\evid{E3} for the PRESENT S-box, sub-linear
  scaling in the number of tables, and a parameter set where the noise price becomes visible.
\end{itemize}
We make no novelty claim: everything here is known since~\cite{CIM19}.
